%% supplemental.tex — Supplemental Material

\paragraph{Fig.~S1 — Subsystem-size dependence.}
The robustness of Z2 QME under variation of the subsystem size $\NA$ is demonstrated
at $N=18$, $\gf=2.0$ for four $\theta$-pairs and $\NA \in \{2,3,4,5,6\}$
(20 combinations total). QME is present in all 20 cases; $\tM$ varies with $\NA$
as expected from the scaling of the initial EA gap [Eq.~(7) of the main text].

\paragraph{Fig.~S2 — QME and DQPT occurrence vs.\ $\gf$.}
Panel (a): crossing time $\tM$ versus post-quench field $\gf$ for the
$\theta$-pair $(0.4,\,1.2)$, distinguishing the FM phase ($\gf < \gc = 0.5$,
blue circles) from the PM phase ($\gf > \gc$, red squares).
Panel (b): binary occurrence of QME and DQPT as a function of $\gf$
($N=14$, $\NA=4$). QME (blue) is present at every tested field; DQPT (orange)
is present only for $\gf > \gc = 0.5$.

\paragraph{Fig.~S3 — XY chain anisotropy dependence.}
Crossing times $\tM$ for the anisotropic XY chain at five values
$\gamma \in \{0, 0.25, 0.5, 0.75, 1.0\}$ for both $\gf=2.0$ (PM phase,
$\gf > \gc = 0.5$) and $\gf=0.42$ (FM phase, $\gf < \gc = 0.5$ for all
$\gamma$; see Eq.~(3) of the main text). QME is present for all five
anisotropies at both fields.

\paragraph{Fig.~S4 — Mode occupancy.}
Quasiparticle mode-occupancy profiles $\langle n_k\rangle(t)$ for selected
$\theta$-pairs, illustrating the differential mode excitation that underlies
the EA crossing dynamics.
